\documentclass[a4paper,14pt]{extarticle}
\usepackage[T2A]{fontenc}
\usepackage[utf8]{inputenc}
\usepackage[russian]{babel}
\usepackage{newtxtext,newtxmath}
\usepackage{geometry}
\usepackage{setspace}
\usepackage{indentfirst}
\usepackage{titlesec}
\usepackage{enumitem}
\usepackage{amsmath}
\usepackage{booktabs}
\usepackage{longtable}
\usepackage{array}
\usepackage{graphicx}
\usepackage{float}
\usepackage{xcolor}
\usepackage{listings}
\usepackage{hyperref}
\usepackage{pgfgantt}

\geometry{left=30mm,right=15mm,top=20mm,bottom=20mm}
\onehalfspacing
\setlength{\parindent}{1.25cm}
\setlist{nosep,leftmargin=1.25cm}

\titleformat{\section}{\normalfont\Large\bfseries}{\thesection}{1em}{}
\titleformat{\subsection}{\normalfont\large\bfseries}{\thesubsection}{1em}{}
\titlespacing*{\section}{0pt}{12pt}{6pt}
\titlespacing*{\subsection}{0pt}{10pt}{4pt}

\hypersetup{colorlinks=true,linkcolor=blue,urlcolor=blue}

\lstset{
  language=Haskell,
  basicstyle=\ttfamily\small,
  keywordstyle=\color{blue}\bfseries,
  commentstyle=\color{green!50!black},
  stringstyle=\color{red!70!black},
  numbers=left,
  numberstyle=\tiny\color{gray},
  stepnumber=1,
  numbersep=5pt,
  showstringspaces=false,
  breaklines=true,
  frame=single
}

\begin{document}

\begin{titlepage}
\thispagestyle{empty}
\begin{center}
{\large МИНИСТЕРСТВО НАУКИ И ВЫСШЕГО ОБРАЗОВАНИЯ РОССИЙСКОЙ ФЕДЕРАЦИИ}\\
{\large «Санкт-Петербургский политехнический университет Петра Великого»}\\
Институт компьютерных наук и кибербезопасности\\
Высшая школа технологий искусственного интеллекта\\
Направление: 02.03.01 «Математика и компьютерные науки»

\vspace{3cm}
{\large Отчет о выполнении курсовой работы}\\
По дисциплине: «Функциональное программирование»\\
На тему: «Учет продуктивности работы сотрудников и расчет заработной платы на основе KPI»
\end{center}

\vspace{3cm}
\noindent Выполнил студент группы 5130201/40002 \hfill \underline{\hspace{3cm}} Храмов В.\,Д.\\[0.8cm]
\noindent Проверил к.\,т.\,н. \hfill \underline{\hspace{3cm}} Моторин Д.\,Е.

\vfill
\begin{center}
Санкт-Петербург\\
2026
\end{center}
\end{titlepage}

\pagenumbering{arabic}
\setcounter{page}{2}

\tableofcontents
\newpage

\section*{Введение}
\addcontentsline{toc}{section}{Введение}

Предметная область работы --- автоматизация учета производственной деятельности сотрудников и расчет заработной платы на основании KPI (Key Performance Indicators). В прикладных информационных системах этот класс задач включает: сбор первичных данных о выполненных задачах, нормализацию и валидацию данных, расчет интегральных показателей эффективности, преобразование KPI в коэффициенты оплаты труда, а также формирование отчетности для пользователя и/или бухгалтерии.

В современном управлении персоналом KPI рассматривается как инструмент количественной оценки результата труда, позволяющий повысить прозрачность мотивации и связать индивидуальный вклад сотрудника с системой вознаграждения. При этом важна не только сама формула KPI, но и инженерная реализация: корректная обработка неполных периодов, устойчивость к ошибочным данным, воспроизводимость расчетов и проверяемость вычислительной логики тестами.

Для реализации курсового проекта была выбрана функциональная парадигма и язык Haskell. Выбор обусловлен строгой системой типов, удобством композиции чистых функций и возможностью изолировать ввод/вывод от доменной логики. Это упрощает формальное обоснование корректности и тестирование (в том числе property-based подходом через QuickCheck).

В работе использованы учебные и прикладные источники по KPI, функциональному программированию и инженерным практикам:
\begin{enumerate}
  \item Marr B. \textit{Key Performance Indicators: The 75 measures every manager needs to know}. Pearson, 2012.
  \item Parmenter D. \textit{Key Performance Indicators: Developing, Implementing, and Using Winning KPIs}. Wiley, 2019.
  \item Bird R. \textit{Thinking Functionally with Haskell}. Cambridge University Press, 2014.
  \item O'Sullivan B., Goerzen J., Stewart D. \textit{Real World Haskell}. O'Reilly, 2008.
  \item Claessen K., Hughes J. QuickCheck: A Lightweight Tool for Random Testing of Haskell Programs // ICFP, 2000.
  \item Документация GHC: \url{https://www.haskell.org/ghc/}.
  \item Документация Stack: \url{https://docs.haskellstack.org/}.
  \item Документация библиотек cassava и aeson: \url{https://hackage.haskell.org/}.
\end{enumerate}

\section{Описание разработки проекта}

\subsection{Архитектура и модульная структура}
Проект реализован в виде консольного приложения с разделением на логические модули:
\begin{itemize}
  \item \texttt{Core} --- доменные типы и общие структуры ошибок.
  \item \texttt{DataLoader} --- загрузка CSV/JSON, парсинг и сериализация выходных отчетов.
  \item \texttt{Productivity} --- вычисление индивидуальной и интегральной продуктивности.
  \item \texttt{Salary} --- вычисление заработной платы с учетом KPI, бонусов и удержаний.
  \item \texttt{Logging} --- журналирование действий пользователя.
  \item \texttt{Utils} --- вспомогательная валидация и работа с датами.
  \item \texttt{app/Main.hs} --- консольный интерфейс и сценарий взаимодействия.
\end{itemize}

\subsection{Технические особенности реализации}
Ключевые технические решения:
\begin{enumerate}
  \item \textbf{Строгая типизация предметной области.} Для сущностей сотрудника, задачи, конфигурации и строк отчетов введены отдельные типы, что снижает риск логических ошибок.
  \item \textbf{Разделение доменной логики и I/O.} Расчет KPI и зарплаты реализован чистыми функциями; операции чтения/записи изолированы в модуле загрузки.
  \item \textbf{Валидация входных данных.} Реализована проверка формата месяца и допустимого диапазона (\texttt{2025-01}...\texttt{2025-12}), обработка отсутствующих файлов и ошибок парсинга.
  \item \textbf{Поддержка неполного месяца.} В конфигурацию введен параметр \texttt{minMonthlyHours}. Для месяцев с меньшей выработкой отмечается статус ``неполный месяц'', а зарплатная формула учитывает покрытие нормы часов.
  \item \textbf{Тестирование.} Использованы property-based тесты QuickCheck для проверок неотрицательности продуктивности и зарплатных показателей.
\end{enumerate}

\subsection{Математическая модель}
Для задачи сотрудника используется обобщенный балл продуктивности:
\[
P_{task} = C \cdot K_c \cdot K_t \cdot K_r,
\]
где \(C\) --- сложность задачи, \(K_c\) --- коэффициент типа/сложности, \(K_t\) --- коэффициент по времени выполнения, \(K_r\) --- коэффициент результата (\texttt{success}/\texttt{partial}/\texttt{failed}).

Интегральный KPI сотрудника за месяц:
\[
KPI = \frac{\sum P_{task}}{N},
\]
где \(N\) --- норматив продуктивности сотрудника.

Итоговая зарплата строится на базе фактически отработанных часов, поправочного коэффициента производительности, бонусной шкалы, а также налоговых и социальных удержаний.

\section{Примеры работы программы}

\subsection{Запуск приложения}
\begin{lstlisting}[language=bash,caption={Запуск проекта}]
cd /Users/vladislav/Documents/fp-cw/kpi-payroll
stack run
\end{lstlisting}

\subsection{Расчет зарплаты за месяц}
\begin{lstlisting}[caption={Фрагмент диалога в консоли}]
=== Учет KPI и расчет зарплаты ===
1. Рассчитать зарплату за месяц (ГГГГ-ММ)
2. Показать топ сотрудников по продуктивности (ГГГГ-ММ)
3. Выгрузить зарплату в CSV (ГГГГ-ММ)
0. Выход
> 1
> 2025-04
\end{lstlisting}

Результат выводится в табличном виде с колонками: ``Сотрудник'', ``KPI'', ``Часы'', ``Статус месяца'', ``К выплате, руб.''.

\subsection{Выгрузка в CSV}
После выбора пункта экспорта формируется файл вида:
\begin{lstlisting}[language=bash,caption={Пример выходного файла}]
output/payroll-2025-04.csv
\end{lstlisting}

CSV содержит расчетные поля, включая новый атрибут \texttt{month\_status} (полный/неполный месяц).

\section{Руководство по эксплуатации для пользователя}

\subsection{Требования}
\begin{itemize}
  \item Установленный GHC и Stack.
  \item Наличие входных файлов в директории \texttt{data/}:
  \texttt{employees.csv}, \texttt{tasks.csv}, \texttt{calendar.txt}, \texttt{config.json}.
\end{itemize}

\subsection{Порядок работы}
\begin{enumerate}
  \item Открыть терминал в каталоге \texttt{kpi-payroll}.
  \item Выполнить \texttt{stack run}.
  \item Выбрать действие в меню:
  \begin{itemize}
    \item \texttt{1} --- расчет зарплаты;
    \item \texttt{2} --- топ сотрудников по продуктивности;
    \item \texttt{3} --- экспорт CSV.
  \end{itemize}
  \item Ввести месяц в формате \texttt{ГГГГ-ММ} в диапазоне \texttt{2025-01}...\texttt{2025-12}.
  \item Для выхода выбрать \texttt{0}.
\end{enumerate}

\subsection{Обработка ошибок}
При некорректном месяце программа не возвращает пользователя в главное меню, а запрашивает повторный ввод месяца. Ошибки и действия пользователя фиксируются в файле \texttt{logs/app.log}.

\section*{Заключение}
\addcontentsline{toc}{section}{Заключение}

В ходе курсовой работы разработан программный комплекс для учета продуктивности сотрудников и расчета заработной платы на основе KPI. Реализованы загрузка и валидация входных данных, расчетные модули, консольный интерфейс, экспорт отчетов и логирование. Программа протестирована и демонстрирует корректную работу на расширенном наборе данных, приближенном к условиям реальной компании.

По результатам разработки выполнены ключевые требования ТЗ: модульность, обработка ошибок, наличие тестов, интерактивное меню, экспорт результатов и воспроизводимость расчетов.

\subsection*{Выгрузка коммитов по этапам}
\addcontentsline{toc}{subsection}{Выгрузка коммитов по этапам}

\begin{longtable}{>{\raggedright\arraybackslash}p{2.6cm} >{\raggedright\arraybackslash}p{3.0cm} >{\raggedright\arraybackslash}p{1.9cm} >{\raggedright\arraybackslash}p{7.0cm}}
\toprule
Этап & Фактический диапазон & Коммит & Сообщение \\
\midrule
\endfirsthead
\toprule
Этап & Фактический диапазон & Коммит & Сообщение \\
\midrule
\endhead
ТЗ & 20.03--10.04 & \texttt{0aa01a4} & Добавлено техническое задание \\
ТЗ & 20.03--10.04 & \texttt{cda53f2} & Исправлено техническое задание \\
ТЗ & 20.03--10.04 & \texttt{5f0a597} & Внесены изменения в техническое задание \\
Диаграммы/структура & 03.04--24.04 & \texttt{b5dd25a} & Сделана диаграмма модулей \\
Диаграммы/структура & 03.04--24.04 & \texttt{05110f0} & Создана структура проекта \\
Реализация & 24.04--06.05 & \texttt{70a54ee} & Сделана минимальная реализация проекта \\
Реализация & 01.05--06.05 & \texttt{123c747} & Валидация по 2025 году, данные \\
Реализация & 01.05--06.05 & \texttt{80b9db2} & Исправлен ввод, фильтрация сотрудников \\
Тестирование/полировка & 06.05 & \texttt{30acecb} & Протестирована работа программы \\
\bottomrule
\end{longtable}

\subsection*{Фактические сроки на диаграмме Ганта}
\addcontentsline{toc}{subsection}{Фактические сроки на диаграмме Ганта}

\begin{figure}[H]
\centering
\begin{ganttchart}[
  hgrid,
  vgrid,
  time slot format=isodate-yearmonthday,
  x unit=0.5cm,
  y unit chart=0.7cm,
  title height=1,
  bar/.append style={fill=gray!30}
]{2026-03-20}{2026-05-22}
\gantttitlecalendar{year, month=name}\\
\ganttbar{Техническое задание (факт)}{2026-03-20}{2026-04-10}\\
\ganttbar{Диаграммы и структура (факт)}{2026-04-03}{2026-04-24}\\
\ganttbar{Реализация программы (факт)}{2026-04-24}{2026-05-06}\\
\ganttbar{Отчет и тестирование (факт)}{2026-04-08}{2026-05-07}
\end{ganttchart}
\caption{Фактические сроки выполнения этапов по данным git-истории}
\end{figure}

\section*{Список литературы}
\addcontentsline{toc}{section}{Список литературы}
\begin{enumerate}
  \item Marr B. Key Performance Indicators: The 75 measures every manager needs to know. Pearson, 2012.
  \item Parmenter D. Key Performance Indicators: Developing, Implementing, and Using Winning KPIs. Wiley, 2019.
  \item Bird R. Thinking Functionally with Haskell. Cambridge University Press, 2014.
  \item O'Sullivan B., Goerzen J., Stewart D. Real World Haskell. O'Reilly, 2008.
  \item Claessen K., Hughes J. QuickCheck: A Lightweight Tool for Random Testing of Haskell Programs // ICFP, 2000.
  \item GHC User's Guide [Электронный ресурс]. URL: \url{https://www.haskell.org/ghc/} (дата обращения: 07.05.2026).
  \item Stack Documentation [Электронный ресурс]. URL: \url{https://docs.haskellstack.org/} (дата обращения: 07.05.2026).
  \item Hackage: cassava, aeson, QuickCheck [Электронный ресурс]. URL: \url{https://hackage.haskell.org/} (дата обращения: 07.05.2026).
\end{enumerate}

\appendix
\section{Приложения}

\subsection{Фрагмент структуры проекта}
\begin{lstlisting}[language=bash]
kpi-payroll/
  app/Main.hs
  data/*.csv, config.json, calendar.txt
  src/Core/*
  src/DataLoader/*
  src/Productivity/*
  src/Salary/*
  src/Logging/*
  src/Utils/*
\end{lstlisting}

\subsection{Формат входного CSV задач}
\begin{lstlisting}
employee_id,date,task_id,complexity,planned_hours,actual_hours,result
E001,2025-04-29,SPRINT_DELIVERY,8,168,163,success
\end{lstlisting}

\end{document}
